\documentclass[11pt]{amsart}
\usepackage{geometry}                % See geometry.pdf to learn the layout options. There are lots.
\geometry{letterpaper}                   % ... or a4paper or a5paper or ... 
%\geometry{landscape}                % Activate for for rotated page geometry
%\usepackage[parfill]{parskip}    % Activate to begin paragraphs with an empty line rather than an indent
\usepackage{graphicx}
\usepackage{amssymb}
\usepackage{epstopdf}
\DeclareGraphicsRule{.tif}{png}{.png}{`convert #1 `dirname #1`/`basename #1 .tif`.png}

\title{Problema 4.1}
\author{Enrique Zuñiga Zuñiga}
%\date{}                                           % Activate to display a given date or no date

\begin{document}
\maketitle
%\section{}
\noindent Escriba un diagrama de flujo que reciba como entrada un arreglo unidimensional ordenado de enteros (posiblemente repetidos) y genere como salida una lista de los números enteros, pero sin repeticiones.  
Dato: VEC [1..N] 1$ <$ N $< $500  \newline
Donde:VEC es un arreglo unidimensional de enteros, cuya capacidad máxima es de 500 elementos.  \newline 
\textbf{Explicación de variables} \newline
N: Variable de tipo entero. Indica el número de elementos que se recibirán como datos. Su valor debe estar comprendido en el intervalo [1..500]. \newline
I: Variable de tipo entero. Se usa como índice del arreglo y como variable de control en diferentes ciclos.
VEC: Arreglo unidimensional de tipo entero. \newline
REPET: Variable de tipo entero. Se utiliza como auxiliar para almacenar el valor impreso. Con este valor se comparan los siguientes números \newline
para determinar si hay repetidos. \newline
A continuación presentamos el programa correspondiente. \newline

%\subsection{}



\end{document}  